\chapter{Damas y ajedrez paso a paso}

\section{Por qué estos juegos son distintos de tres en raya}

Tres en raya permite mostrar el árbol completo de posiciones pequeñas. Damas y ajedrez introducen dos dificultades nuevas:

\begin{enumerate}
\item el factor de ramificación es mucho mayor;
\item las hojas evaluadas normalmente no son terminales reales, sino posiciones cortadas a cierta profundidad.
\end{enumerate}

Por eso, en estos juegos distinguimos entre terminalidad real y terminalidad por profundidad.

\begin{definition}[Hoja artificial]
Una hoja artificial es un estado no terminal que el análisis decide tratar como hoja porque se alcanzó una profundidad máxima de búsqueda.
\end{definition}

\section{Damas: ejemplo de captura obligatoria}

Consideremos una posición simplificada de damas. Usaremos $x$ para una ficha de MAX, $o$ para una ficha de MIN y $\blank$ para una casilla vacía relevante.

\[
D_0=
\begin{array}{cccc}
\blank & o & \blank & \blank\\
\blank & \blank & x & \blank\\
\blank & \blank & \blank & \blank\\
\blank & \blank & \blank & \blank
\end{array}
\]

Si las reglas obligan a capturar cuando existe una captura, el conjunto de movimientos legales no contiene todos los movimientos diagonales posibles, sino solo los movimientos de captura.

\begin{formalbox}
La diferencia se expresa así:
\[
\Legal(D_0)=\{a\in\Moves: a\text{ es captura legal desde }D_0\},
\]
si existe al menos una captura. Solo cuando no existe captura se permiten movimientos ordinarios.
\end{formalbox}

\begin{intuitionbox}
La regla de captura obligatoria reduce el árbol. En vez de ramificar sobre muchos movimientos tranquilos, el juego fuerza una línea táctica. Esto puede hacer que Alpha--Beta encuentre refutaciones más rápido.
\end{intuitionbox}

\section{Árbol inicial de damas}

\[
\begin{forest}
game tree
[$D_0$\\MAX mueve
  [$D_1$\\captura obligatoria]
]
\end{forest}
\]

Este árbol tiene un solo hijo en el primer nivel. Por tanto,
\[
 b(D_0)=1.
\]
Aunque el juego de damas suele tener muchas posibilidades, las reglas locales pueden hacer que el factor de ramificación cambie drásticamente de un nodo a otro.

\section{Ajedrez: ejemplo de apertura}

En ajedrez, desde la posición inicial las blancas tienen 20 movimientos legales. Esto ya produce un árbol mucho más grande que tres en raya.

Representaremos una línea inicial:
\[
1.e4\quad c5\quad 2.Nf3.
\]

\begin{stepbox}{Paso 1: raíz}
La raíz $C_0$ es la posición inicial. Mueven blancas, que tomaremos como MAX.
\end{stepbox}

\begin{stepbox}{Paso 2: primer movimiento}
MAX considera $1.e4$. El estado sucesor se llama $C_1$.
\end{stepbox}

\begin{stepbox}{Paso 3: respuesta de MIN}
MIN responde $1...c5$. El estado sucesor se llama $C_2$.
\end{stepbox}

\begin{stepbox}{Paso 4: segundo movimiento de MAX}
MAX juega $2.Nf3$. El estado sucesor se llama $C_3$.
\end{stepbox}

\section{Árbol parcial de ajedrez}

\[
\begin{forest}
game tree
[$C_0$\\blancas/MAX
  [$C_1$\\$1.e4$
    [$C_2$\\$1...c5$
      [$C_3$\\$2.Nf3$]
      [$C_4$\\$2...d6$]
      [$C_5$\\$2...Nc6$]
    ]
    [$C_6$\\$1...e5$]
    [$C_7$\\$1...c6$]
  ]
  [$C_8$\\$1.d4$]
  [$C_9$\\$1.Nf3$]
]
\end{forest}
\]

\section{Evaluación de hojas artificiales en ajedrez}

En ajedrez real no podemos llegar al final de la partida. Por eso, al cortar el árbol, usamos una función de evaluación heurística.

\begin{definition}[Evaluación heurística]
Una evaluación heurística es una función
\[
E:\States\to\mathbb{R}
\]
que aproxima la utilidad de una posición no terminal.
\end{definition}

Una evaluación básica puede combinar:
\begin{itemize}
\item material;
\item movilidad;
\item seguridad del rey;
\item control del centro;
\item estructura de peones.
\end{itemize}

\begin{warningbox}
La evaluación heurística no es igual a la utilidad terminal. Minimax con hojas artificiales calcula el valor óptimo respecto a la evaluación elegida, no necesariamente respecto al resultado real de la partida.
\end{warningbox}

\section{Plan de ampliación}

Las siguientes entregas convertirán este bosquejo en ejemplos completos:

\begin{enumerate}
\item diagramas de tablero para cada nodo;
\item explicación de cada movimiento legal considerado;
\item valores heurísticos de hojas;
\item retropropagación Minimax;
\item poda Alpha--Beta con $\alpha$ y $\beta$ en cada nodo;
\item comparación entre orden bueno y orden malo de movimientos.
\end{enumerate}
